\section{Collective-Intelligence Prevention Relevance}
\label{sec:prevention-relevance}

The \texttt{prevention\_relevance} module evaluates whether historical threat
intelligence is applicable to a proposed organizational change. It represents
sample system and network telemetry for the proposed deployment as a typed
graph $G_{\mathrm{new}}$, converts historical STIX Indicators and their context
into candidate graphs, and ranks those Indicators using an interpretable graph
similarity score. The output supports analyst review; relevance does not prove
that an environment is vulnerable, compromised, or malicious.

\subsection{Inputs and Processing}

The proposed-change input consists of one or more CSV files containing
representative system or network events. The module reuses the field-aware
preprocessing functions of the federated pipeline, specifically
\texttt{field\_aware\_tokens()} and \texttt{\_row\_text()} from
\texttt{federated\_lr\_pipeline.specialized\_models}. It does not fit a model,
build a vocabulary, or invoke federated inference.

Each resolved value becomes a typed graph node, for example a platform,
process, file path, IP address, domain, protocol, or port. Relations encode
observed structure, including \texttt{runs}, \texttt{accesses},
\texttt{connects\_to}, \texttt{uses\_protocol}, and \texttt{uses\_port}.
Canonical values produce deterministic node and edge identifiers. CSV aliases
are defined in YAML/JSON configuration; conflicting populated aliases are
omitted conservatively.

The historical input may be a STIX Bundle, a JSON list/object, or JSONL STIX
objects. Exact STIX equality patterns for IP addresses, domains, URLs, email
addresses, vulnerabilities, and common file hashes are converted to observable
nodes. Optional IoC records and model explanations link Indicators to attack
classes, source events, evidence categories, and field-aware features.

\subsection{Typed Similarity}

Let $C_j$ be the ego graph of Indicator $j$, let $N$ be the number of candidate
Indicators, and let $\operatorname{df}(u)$ be the number of candidate graphs
containing node $u$. A candidate node receives weight
\begin{equation}
 w(u)=w_{\operatorname{type}(u)}
 \left[
 \log\left(\frac{N+1}{\operatorname{df}(u)+1}\right)+1
 \right].
\end{equation}
This reduces the influence of evidence shared by many candidates.

The node-coverage component is
\begin{equation}
 S_{\mathrm{node}}(C_j,G_{\mathrm{new}})=
 \frac{
 \sum_{u\in C_j}w(u)\max_{v\in G_{\mathrm{new}}}
 \operatorname{sim}(u,v)
 }{
 \sum_{u\in C_j}w(u)
 }.
\end{equation}
The typed similarity function supports exact canonical equality,
IP-in-network membership, process/path basename agreement, and URL-host/domain
agreement. The edge component $S_{\mathrm{edge}}$ applies an analogous weighted
coverage calculation to relation-equivalent edges, using the product of the
two endpoint similarities. $S_{\mathrm{direct}}$ equals one when the exact
observable named by the Indicator occurs in $G_{\mathrm{new}}$; otherwise it is
zero. $S_{\mathrm{context}}$ measures configured attack-class applicability,
such as platform, process, protocol, or port requirements.

The final score is
\begin{equation}
 S_j=\operatorname{clip}_{[0,1]}\left(
 \alpha_dS_{\mathrm{direct}}+
 \alpha_nS_{\mathrm{node}}+
 \alpha_eS_{\mathrm{edge}}+
 \alpha_cS_{\mathrm{context}}-
 \lambda I_{\mathrm{platform}}
 \right),
\end{equation}
where $I_{\mathrm{platform}}$ identifies an explicit platform incompatibility
and the $\alpha$ values are normalized configuration weights. Unless an exact
direct observable matches, a candidate must satisfy a distinctive-evidence
gate and have edge or contextual support. Consequently, generic evidence such
as TCP or port 443 cannot select an Indicator by itself.

Scores above the high threshold are marked \texttt{relevant}; scores in the
lower review band are marked \texttt{analyst\_review}; all others are
\texttt{not\_selected}. A separate priority score incorporates Indicator
validation status and an available model risk probability for ordering, but it
does not change the architectural relevance score.

\subsection{Execution}

\begin{verbatim}
conda run -n LR python -m prevention_relevance.run \
  --proposed-system-logs proposed/system.csv \
  --proposed-network-logs proposed/network.csv \
  --ioc-bundle outputs/run_001/iocs/ioc_bundle.json \
  --ioc-records outputs/run_001/iocs/ioc_records.jsonl \
  --explanations outputs/run_001/explanations.jsonl \
  --config prevention_relevance/examples/similarity.yaml \
  --output-dir outputs/prevention-run
\end{verbatim}

The command writes $G_{\mathrm{new}}$, the historical graph, CSV and JSONL
rankings, per-Indicator explanations, a run summary, and the resolved
configuration. Historical-graph serialization may be disabled for large
bundles. If an analyst label file with columns \texttt{indicator\_id} and
\texttt{relevant} is supplied, the module additionally reports classification
and ranking metrics, including precision, recall, F1, average precision, ROC
AUC, NDCG, and precision/recall at $k$.

\subsection{Interpretation and Limitations}

The analysis is deterministic and performs no DNS resolution, online lookup,
fuzzy matching, or graph embedding. Its quality depends on the completeness of
the proposed-change telemetry, the STIX context, and administrator-defined
applicability mappings. An exact overlap or high similarity indicates that
historical intelligence should be reviewed when planning monitoring or
controls; it is not evidence that the proposed deployment is benign or
malicious. Unreviewed IoCs and administrative applicability assumptions must
remain subject to human validation.

